\scanpage{144}
\margin{\textcircled{32}}
\uline{Formule de Lefschetz-Verdier.}
\diagram{\begin{tikzcd}[column sep=small,row sep=small]
&X\times_S Y\arrow[dl,no head,"\operatorname{pr}_1"']\arrow[dr,no head,"\operatorname{pr}_2"]\arrow[dd,"h" near start]&\\
F_\bullet\quad X\arrow[dr,no head,"f"']&&Y\quad G_\bullet\arrow[dl,no head,"g"]\\
&S&
\end{tikzcd}}
\margin{\textcircled{a}}
\[
\begin{array}{l}
\operatorname{Hom}^0(\operatorname{pr}_1^*(F_\bullet),\operatorname{pr}_2^!(G_\bullet))\simeq\mathbb R\Gamma_{X\times_S Y}\mathbb R\mathcal Hom^\bullet(\operatorname{pr}_1^*(F_\bullet),\operatorname{pr}_2^!(G_\bullet))\\
\quad\rotatebox{90}{$\simeq$}\\
\operatorname{Hom}^0(\operatorname{pr}_{2!}\operatorname{pr}_1^*(F_\bullet),G_\bullet)\\
\quad\rotatebox{90}{$\simeq$}\\
\operatorname{Hom}^0(g^*f_!(F_\bullet),G_\bullet)\\
\quad\rotatebox{90}{$\simeq$}\\
\operatorname{Hom}^0(f_!(F_\bullet),g_*(G_\bullet))
\end{array}
\]
\margin{correspondances cohomologiques de degré $0$ de $(X,F_\bullet)$ dans $(Y,G_\bullet)$ sur $S$ \dots}
\[
\operatorname{Hom}^0(\operatorname{pr}_1^*(F_\bullet),\operatorname{pr}_2^!(G_\bullet))\simeq\operatorname{Hom}^0(f_!(F_\bullet),g_*(G_\bullet))
\]
\[
\operatorname{Hom}^0(\operatorname{pr}_2^*(G_\bullet),\operatorname{pr}_1^!(F_\bullet))\simeq\operatorname{Hom}^0(g_!(G_\bullet),f_*(F_\bullet))
\]
Supposons que $f$ et $g$ soient propres, alors les
deux expressions précédentes s'accouplent dans
\deleted{\unsure{$\operatorname{Hom}^0(Rf_!$}}
\[
\operatorname{Hom}^0(f_*(F_\bullet),\unsure{f_*|f}(F_\bullet))\quad\text{et dans}\quad\operatorname{Hom}^0(g_*(G_\bullet),g_*(G_\bullet))
\]
trace, dans \deleted{\unsure{$\operatorname{Hom}$}}\deleted{\unsure{$A$}} $A_S$
\inserted{si $S$ discret, et $F_\bullet$ et $G_\bullet$ de Tor-dim finie}
(par ``sommes alternées
des traces'').
\deleted{Cet accouplement est induit
par un accouplement de nature locale (\unsure{pas besoin}
que $f$ et $g$ soient propres)
\[
\mathbb R\mathcal Hom^\bullet(\operatorname{pr}_1^*(F_\bullet),\operatorname{pr}_2^!(G_\bullet))\times\mathbb R\mathcal Hom(\operatorname{pr}_2^*(G_\bullet),\operatorname{pr}_1^!(F_\bullet))\longrightarrow h^!(A_S)
\]
\unsure{$=$}
qu'on définit ainsi}
\figure{\begin{tikzpicture}\draw (2,0)--(9,0);\draw (0,-.7)--(2,-.7);\draw (0,-.7)--(0,-2.4)--(4.1,-2.4)--(4.3,-1.6)--(9,-1.6);\draw (.8,-1)--(2,-2.2);\draw (3,-.4)--(4,-1.4);\draw (6,0)--(7,-1.3);\draw (8.5,0)--(9,-1);\draw (4,-2.8)--(5.5,-2.8);\draw (4.4,-2.9)--(5.4,-2.9);\draw (4.8,-3)--(5.3,-3);\end{tikzpicture}}
\mspage{143}
